% !TEX root = chapter-standalone.tex

\section{Traced symmetric strict monoidal categories}


\todotext{make a remark that without symmetry we would have two distinct concepts of left-trace and right-trace}


\begin{ctdefinition}[Traced symmetric strict monoidal category]
    \label{def:traced-strict-monoidal-cat}
    \label{def:traced-strict-monoidal-category}
    We say that a \SY{symmetric strict monoidal category}~$\tup{\CatC, \mtimescat, \idmoncat, \braiding}$ is \emph{traced} if it is equipped with a \emph{trace operator}: a family of functions
    \begin{equation}
        \Tr_{\Obja,\Objb}^\Objc\colon \HomSet{\CatC}{\Obja \mtimescatob \Objc}{\Objb\mtimescatob \Objc}\to \HomSet{\CatC}{\Obja}{\Objb},
    \end{equation}
    satisfying the following axioms:
    \begin{enumerate}
        \item \emph{Naturality in $\Obja$:} For any morphisms $\mora\colon \Obja\mtimescatob \Objc \mto \Objb\mtimescatob \Objc$ and $\morb \colon \Obja' \mto \Obja$,
              \begin{equation}\label{eq:trace-naturality-1}
                  \Tr_{\Obja', \Objb}^\Objc ( (\morb \mtimescatmor \catidat\Objc) \mthen \mora) = \morb \mthen \Tr_{\Obja, \Objb}^\Objc (\mora)
              \end{equation}
        \item \emph{Naturality in $\Objb$:} For any morphisms $\mora\colon \Obja\mtimescatob \Objc \mto \Objb\mtimescatob \Objc$ and $\morb \colon \Objb \mto \Objb'$,
              \begin{equation}\label{eq:trace-naturality-2}
                  \Tr_{\Obja, \Objb'}^\Objc ( \mora \mthen (\morb \mtimescatmor \catidat\Objc) ) = \Tr_{\Obja, \Objb}^\Objc (\mora) \mthen \morb
              \end{equation}
        \item \emph{Dinaturality in $\Objc$:}
              For any morphisms $\mora \colon \Obja  \mtimescatob \Objc \mto \Objb \mtimescatob \Objc'$ and $\morb \colon \Objc' \mto \Objc$,
              \begin{equation}\label{eq:trace-dinaturality}
                  \Tr_{\Obja, \Objb}^{\Objc} (\mora \mthen (\catidat\Objb \mtimescatmor \morb)) = \Tr_{\Obja, \Objb}^{\Objc'}((\catidat\Obja \mtimescatmor \morb) \mthen \mora).
              \end{equation}

        \item \emph{Vanishing I:}
              For any morphisms~$\mora \colon \Obja \mtimescatob \idmoncat \mto \Objb \mtimescatob \idmoncat$ in \CatC,
              \begin{equation}
                  \label{eq:vanishing_1}
                  \Tr_{\Obja,\Objb}^\idmoncat (\mora)=  \mora.
              \end{equation}
        \item \emph{Vanishing II:}
              For any morphism~$\mora \colon (\Obja\mtimescatob \Objc) \mtimescatob \Objd \mto (\Objb\mtimescatob \Objc) \mtimescatob \Objd$ in \CatC,
              \begin{equation}
                  \label{eq:vanishing_2}
                  \Tr_{\Obja,\Objb}^\Objc\pars{
                      \Tr_{\Obja \mtimescatob \Objc, \Objb \mtimescatob \Objc}^\Objd(\mora)} = \Tr_{\Obja,\Objb}^{\Objc\mtimescatob \Objd}( \mora ).
              \end{equation}
        \item \emph{Superposing:}
              For any morphism~$\mora\colon \Obja\mtimescatob \Objc \mto \Objb\mtimescatob \Objc$ in \CatC,
              \begin{equation}
                  \label{eq:superposing}
                  \Tr_{\Obje\mtimescatob \Obja,\Obje\mtimescatob \Objb}^{\Objc}( \catidat\Obje\mtimescatmor \mora )=\catidat\Obje\mtimescatmor \Tr_{\Obja,\Objb}^\Objc(\mora).
              \end{equation}
        \item \emph{Yanking:}
              \begin{equation}
                  \label{eq:yanking}
                  \Tr_{\Objc,\Objc}^\Objc\pars{\braiding_{\Objc,\Objc}}=\catidat\Objc.
              \end{equation}
    \end{enumerate}
\end{ctdefinition}

\begin{figure}[h!]
    \centering
    \begin{tikzpicture}[scale=0.7, every node/.style={transform shape},
        morbox/.style={draw=morphisms, fill=white, thick, minimum width=1.6cm, minimum height=2.2cm},
        wire/.style={objects, thick},
        wirelabel/.style={objects, font=\bfseries},
        tracebox/.style={draw=black, dashed, thick, fill=graphcolor, rounded corners=2pt},
    ]
        % Trace boundary (drawn first so it's behind everything)
        \fill[graphcolor, rounded corners=2pt] (-0.6,-2.4) rectangle (4.6,1.6);
        \draw[black, dashed, thick, rounded corners=2pt] (-0.6,-2.4) rectangle (4.6,1.6);

        % Morphism box f
        \node[morbox] (f) at (2,0) {};
        \node[morphisms, font=\large\itshape] at (2,0) {$f$};

        % X wire: input from left
        \draw[wire] (-1.8, 0.55) -- (f.west |- 0,0.55);
        \node[wirelabel, above] at (-1.2, 0.55) {$\Obja$};

        % Y wire: output to right
        \draw[wire] (f.east |- 0,0.55) -- (5.8, 0.55);
        \node[wirelabel, above] at (5.2, 0.55) {$\Objb$};

        % Z wire: input side (from feedback loop up to f)
        \draw[wire] (0.6, -1.8) -- (0.6, -0.55) -- (f.west |- 0,-0.55);
        \node[wirelabel, left] at (0.6, -1.1) {$\Objc$};

        % Z wire: output side (from f down to feedback loop)
        \draw[wire] (f.east |- 0,-0.55) -- (3.4, -0.55) -- (3.4, -1.8);
        \node[wirelabel, right] at (3.4, -1.1) {$\Objc$};

        % Feedback loop: bottom horizontal segment
        \draw[wire] (0.6, -1.8) -- (3.4, -1.8);
    \end{tikzpicture}
    \caption{The trace operator: given a morphism~$\mora \colon \Obja \mtimescatob \Objc \mto \Objb \mtimescatob \Objc$, its trace~$\Tr_{\Obja,\Objb}^\Objc(\mora) \colon \Obja \mto \Objb$ is obtained by ``closing the loop'' on the~$\Objc$ wire.}
    \label{fig:trace-operator}
\end{figure}

\begin{figure}[h!]
    \centering
    \includegraphics[width=10cm]{trace-naturality-x}
    \caption{Naturality in $\Obja$}
    \label{fig:trace-naturality-x}
\end{figure}

\begin{figure}[h!]
    \centering
    \includegraphics[width=10cm]{trace-naturality-y}
    \caption{Naturality in $\Objb$}
    \label{fig:trace-naturality-y}
\end{figure}

\begin{figure}[h!]
    \centering
    \includegraphics[width=10cm]{trace-dinaturality-z}
    \caption{Dinaturality in $\Objc$}
    \label{fig:trace-dinaturality-z}
\end{figure}

\begin{figure}[h!]
    \centering
    \includegraphics[width=6cm]{trace-vanishing-i}
    \caption{Vanishing I}
    \label{fig:trace-vanishing-i}
\end{figure}

\begin{figure}[h!]
    \centering
    \includegraphics[width=10cm]{trace-vanishing-ii}
    \caption{Vanishing II}
    \label{fig:trace-vanishing-ii}
\end{figure}

\begin{figure}[h!]
    \centering
    \includegraphics[width=8cm]{trace-yanking}
    \caption{Yanking}
    \label{fig:trace-yanking}
\end{figure}

\todotext{2026-09-01, JL: change superposing axiom in the definition to be the stronger one, and turn the lemma below into a remark.}

\todotext{2026-09-01, JL: the string diagram for the superposing axiom is missing below; needs to be added.}

\begin{remark}
    Other variants of the definition of a \SY{traced monoidal category} can be found in the literature.
    For instance, some include a more general version of the superposing law, see \cref{lem:general-superposing-law} below.
\end{remark}

\begin{lemma}
    \label{lem:general-superposing-law}
    Let~$\tup{\CatC, \mtimescat,\idmoncat_{\CatC}, \braiding, \Tr}$ be a \SY{traced symmetric strict monoidal category}.
    Then a more general version of the superposing law holds: for any morphisms~
    $$\mora\colon \Obja\mtimescatob \Objc \mto \Objb \mtimescatob \Objc \quad \text{ and } \quad \morb \colon \Objd \mto \Obje,$$
    we have 
    \begin{equation}
        \label{eq:superposing_2}
        \Tr_{\Objd\mtimescatob \Obja,\Obje\mtimescatob \Objb}^{\Objc}(\morb \mtimescatmor \mora)=\morb \mtimescatmor \Tr_{\Obja,\Objb}^\Objc(\mora).
    \end{equation}
\end{lemma}

\todotextjira{165}{\alphubel: @JL: Do the proof.
    Also, this could be made into a graded exercise.
}




